\section{Simulation}

The preceding solutions all operate on the infinite game.
This section asks what the divergence looks like in a finite sample, since any observable version of the game, in a laboratory or elsewhere, is finite.

The design is a Monte Carlo experiment.
A single round draws coin flips until the first head appears, so the number of flips $k$ is geometric with parameter $1/2$ and the payoff is $2^{k}$.
The sample mean is driven entirely by the largest realisation of $k$ in the sample: a draw of $k$ contributes $2^{k}/N$ to the mean, while the probability of observing it is $2^{-k}$.
The experiment runs over \simGridPoints{} sample sizes spaced logarithmically from one round to \simNEnd{} rounds.
At each sample size the entire experiment is repeated \simReps{} times independently, which yields the sampling distribution of the mean payoff.
I report the median across replications rather than the mean, since the latter is subject to the same divergence as the object it would summarise.

\begin{figure}[htbp]
    \centering
    \includegraphics[width=0.8\textwidth]{ensemble_means.pdf}
    \caption{Sampling distribution of the mean payoff against sample size, over \simReps{} independent replications. The solid line is the median across replications, the shaded regions are quantile bands. The horizontal axis is logarithmic.}
    \label{fig:ensemble_means}
\end{figure}

Figure~\ref{fig:ensemble_means} shows no convergence.
The median sample mean rises from \simMedianStart{} at \simNStart{} rounds to \simMedianEnd{} at \simNEnd{} rounds. The linearity in log-scale suggest that the median itself diverges logarithmic.
Furthermore, the distribution is sharply asymmetric. The lower quantiles track the median closely, because the small-$k$ draws that make up most of any sample concentrate quickly, whereas the upper band is outlier is wide and does not settle as the number of replications grows.
The upper quantile does not stabilise as the number of replications grows. This is the display of the paradox itself: the sample mean is governed by rare large realisations, so no finite sample, and no finite number of samples, pins it down.